\documentclass[12pt]{article}
 
\usepackage{fancyhdr}
\usepackage[a4paper, total={6in, 8in}]{geometry}
\usepackage{paralist}
\setlength{\headheight}{15pt} 
\pagestyle{fancy}
\fancyhf{}
\rhead{\today}
\lhead{C. Woodgate}
%\fancyfoot[C]{\thepage}

\begin{document}

\begin{center}
	\large{\textbf{The Idiot's Guide to Fletching}}
\end{center}

These instructions describe how to fletch an aluminium or carbon arrow with a plastic vane, using conventional fletching glue. For wooden arrows, `spin vanes', and feathers you will need to use different techniques.

\subsubsection*{The Ground Rules}

\begin{itemize}
	\item {\bf NEVER} stick a used fletching back on an arrow unless you absolutely have to. Used fletchings have residual glue, grease, and dirt on them and will never stick as well as a fresh one.
	\item {\bf ALWAYS} prepare thoroughly. Set aside enough time and make sure you have the right equipment and materials. Nine times out of ten, when fletchings fall off it is because the arrow and fletching were not prepared correctly.
\end{itemize}

\subsubsection*{What You'll Need}
\begin{compactitem}
	\item Fletching Jig.
	\item Fletching Glue.
	\item New Fletching.
	\item Craft knife.
	\item Kitchen Roll, tissue, or cotton pads.
	\item Solvent: Acetone is best. Methylated spirits or rubbing alcohol are also fine. {\bf AVOID} white spirit and nail polish remover. Both of these tend to leave an oily residue which stops the fletching sticking.
\end{compactitem}

\subsubsection*{How to Prepare your Arrow and Vane}
\begin{enumerate}
	\item {\bf WASH YOUR HANDS.} Grease and dirt are the enemy! It also helps to wipe down the table you will be using if it is mucky.
	\item Remove the old fletching. Normally it is sufficient to pull it off gently with your fingers, but you can also get special `fletching remover' tools.
	\item Scrape any residual glue from the shaft using the craft knife. Go gently and at a shallow angle to avoid damaging the shaft. Be particularly careful with carbon arrows as they are easy to damage!
	\item Wet the cloth you are using with a little of the solvent and use this to thoroughly clean the bit of the arrow to which you will be sticking fletchings.
	\item Finally, clean the base of the vane with the solvent, too. This ensures there is no residue left from your fingers or when it was manufactured.
\end{enumerate}


\subsubsection*{How to Stick the Fletching on}
\begin{enumerate}
	\item Place the arrow in the fletching jig, making sure it is the right way up. Put the clamp on and check that it sits straight down the middle of the arrow shaft. Note which marking to use to put the fletching in the right place in the clamp.
	\item Remove the clamp, and put the fletching in right up to the base, making sure to line it up with the correct mark.
	\item Run a thin bead of fletching glue down the base of the fletching.
	\item Place the clamp on the jig and gently push the fletching down onto the arrow, making sure it is in the correct position. It is normal to see a little glue squeeze out around the base of the vane.
	\item Wait 10 minutes for the glue to set.
	\item Carefully remove the clamp, then take the arrow out of the jig. The final (optional) step is to `top-and-tail' the fletch. This means putting a small dab of glue at both the top and bottom of the fletching to help hold it in place.
	\item All done! Now leave the arrow overnight so the glue can cure fully before shooting it.
\end{enumerate}

\end{document}
